\chapter{计算几何与计算技巧}\label{ux8ba1ux7b97ux51e0ux4f55ux4e0eux8ba1ux7b97ux6280ux5de7}

\section{浮点规范}\label{ux6d6eux70b9ux89c4ux8303}

整数坐标优先用 \passthrough{\lstinline!long long!} 叉积，乘法可能到 \passthrough{\lstinline!1e18!} 时用 \passthrough{\lstinline!\_\_int128!}。实数几何统一 \passthrough{\lstinline!long double!}，比较用 \passthrough{\lstinline!sgn(x) = (x>eps)-(x<-eps)!}，不要直接 \passthrough{\lstinline!==!}。

\begin{lstlisting}[language={C++}]
using Real = long double;
const Real EPS = 1e-12L;
struct Point { Real x, y;
    Point operator+(Point b) const { return {x + b.x, y + b.y}; }
    Point operator-(Point b) const { return {x - b.x, y - b.y}; }
    Point operator*(Real k) const { return {x * k, y * k}; }
};
Real cross(Point a, Point b) { return a.x * b.y - a.y * b.x; }
Real dot(Point a, Point b) { return a.x * b.x + a.y * b.y; }
int sgn(Real x) { return (x > EPS) - (x < -EPS); }
\end{lstlisting}

\section{向量与叉积}\label{ux5411ux91cfux4e0eux53c9ux79ef}

\passthrough{\lstinline!cross(a,b)=a.x*b.y-a.y*b.x!}：正为逆时针，负为顺时针，0 共线。点 \passthrough{\lstinline!p!} 在有向线段 \passthrough{\lstinline!ab!} 左侧当且仅当 \passthrough{\lstinline!cross(b-a,p-a)>0!}。

点积判断投影/夹角：\passthrough{\lstinline!dot(a,b)>0!} 锐角，\passthrough{\lstinline!<0!} 钝角；距离平方尽量不开放平方根。

\begin{lstlisting}[language={C++}]
Real dist2(Point a, Point b) { Point d = a - b; return dot(d, d); }
bool left_of(Point a, Point b, Point p) { return sgn(cross(b - a, p - a)) > 0; }
\end{lstlisting}

\section{直线、线段与交点}\label{ux76f4ux7ebfux7ebfux6bb5ux4e0eux4ea4ux70b9}

线段相交：先做包围盒快速排除，再判断四个方向叉积；共线时检查投影是否重叠。求交点用参数方程：\passthrough{\lstinline!a + t(b-a)!}，分母是两方向向量叉积；平行/重合要单独处理。

点到直线距离：\passthrough{\lstinline!abs(cross(b-a,p-a))/|b-a|!}；点到线段要先看投影是否落在段内，否则取端点距离。

\begin{lstlisting}[language={C++}]
bool on_segment(Point a, Point b, Point p) {
    return sgn(cross(b - a, p - a)) == 0 && sgn(dot(p - a, p - b)) <= 0;
}
bool intersect(Point a, Point b, Point c, Point d) {
    Real c1 = cross(b - a, c - a), c2 = cross(b - a, d - a);
    Real c3 = cross(d - c, a - c), c4 = cross(d - c, b - c);
    if (sgn(c1) == 0 && on_segment(a, b, c)) return true;
    if (sgn(c2) == 0 && on_segment(a, b, d)) return true;
    return sgn(c1) * sgn(c2) < 0 && sgn(c3) * sgn(c4) < 0;
}
\end{lstlisting}

\subsection{圆几何}\label{ux5706ux51e0ux4f55}

直线与圆交点先投影到直线，再根据圆心到直线的距离判断交点数；两圆交点用圆心连线方向分解。判别相切时仍使用 \passthrough{\lstinline!EPS!}。

\begin{lstlisting}[language={C++}]
struct Circle { Point o; Real r; };
vector<Point> line_circle(Point a, Point b, Circle c) {
    Point d = b - a;
    Real dd = dot(d, d), t = dot(c.o - a, d) / dd;
    Point foot = a + d * t;
    Real h2 = c.r * c.r - dist2(foot, c.o);
    if (h2 < -EPS) return {};
    if (h2 <= EPS) return {foot};
    Point unit = d * (sqrtl(h2 / dd));
    return {foot - unit, foot + unit};
}

vector<Point> circle_circle(Circle a, Circle b) {
    Point d = b.o - a.o; Real dd = dot(d, d), dis = sqrtl(dd);
    if (dis < EPS || dis > a.r + b.r + EPS || dis < fabsl(a.r - b.r) - EPS) return {};
    Real x = (a.r * a.r - b.r * b.r + dd) / (2 * dis);
    Real h2 = max((Real)0, a.r * a.r - x * x);
    Point base = a.o + d * (x / dis);
    Point perp{-d.y / dis, d.x / dis};
    if (h2 <= EPS) return {base};
    Point off = perp * sqrtl(h2);
    return {base - off, base + off};
}
\end{lstlisting}

\section{多边形}\label{ux591aux8fb9ux5f62}

\begin{itemize}
\tightlist
\item
  有向面积：\passthrough{\lstinline!Σ cross(p[i], p[(i+1)\%n]) / 2!}。
\item
  点在多边形：射线法（奇偶交点）或绕数法；边界先单独判定。
\item
  凸多边形内点可二分扇区，查询 \passthrough{\lstinline!O(log n)!}；一般多边形射线法 \passthrough{\lstinline!O(n)!}。
\item
  Pick 定理：格点多边形 \passthrough{\lstinline!A = I + B/2 - 1!}，其中 \passthrough{\lstinline!B=Σgcd(|dx|,|dy|)!}。
\end{itemize}

\begin{lstlisting}[language={C++}]
Real area2(const vector<Point>& p) {
    Real s = 0;
    for (int i = 0; i < (int)p.size(); ++i)
        s += cross(p[i], p[(i + 1) % p.size()]);
    return s;
}
\end{lstlisting}

\section{凸包（Andrew）}\label{ux51f8ux5305andrew}

点按 \passthrough{\lstinline!(x,y)!} 排序去重；维护下凸/上凸，若 \passthrough{\lstinline!cross <= 0!} 弹出可得到不保留共线点的凸包，若要保留边界改为 \passthrough{\lstinline!<0!}。复杂度 \passthrough{\lstinline!O(n log n)!}。周长、面积、最远点对都以凸包为输入。

\begin{lstlisting}[language={C++}]
vector<Point> convex_hull(vector<Point> p) {
    sort(p.begin(), p.end(), [](Point a, Point b) {
        return a.x < b.x || (a.x == b.x && a.y < b.y);
    });
    p.erase(unique(p.begin(), p.end(), [](Point a, Point b) {
        return a.x == b.x && a.y == b.y;
    }), p.end());
    vector<Point> h;
    for (Point x : p) {
        while (h.size() >= 2 && cross(h.back() - h[h.size() - 2], x - h.back()) <= 0) h.pop_back();
        h.push_back(x);
    }
    int lower = h.size();
    for (int i = (int)p.size() - 2; i >= 0; --i) {
        while ((int)h.size() > lower && cross(h.back() - h[h.size() - 2], p[i] - h.back()) <= 0) h.pop_back();
        h.push_back(p[i]);
    }
    if (h.size() > 1) h.pop_back();
    return h;
}
\end{lstlisting}

\section{旋转卡壳}\label{ux65cbux8f6cux5361ux58f3}

在凸包上用一个单调指针寻找对踵点，求直径（最远点对）\passthrough{\lstinline!O(h)!}；也可求最大三角形面积、两凸多边形距离。指针前进依据叉积面积是否增大，不要用浮点角度比较。

\begin{lstlisting}[language={C++}]
long double diameter2(const vector<Point>& p) {
    int n = p.size(), j = 1; long double ans = 0;
    for (int i = 0; i < n; ++i) {
        int ni = (i + 1) % n;
        while (fabsl(cross(p[ni] - p[i], p[(j + 1) % n] - p[i])) >
               fabsl(cross(p[ni] - p[i], p[j] - p[i]))) j = (j + 1) % n;
        ans = max(ans, max(dist2(p[i], p[j]), dist2(p[ni], p[j])));
    }
    return ans;
}
\end{lstlisting}

\section{扫描线}\label{ux626bux63cfux7ebf}

矩形并面积：按 x 端点事件排序，y 区间用线段树维护覆盖长度，当前面积增量为 \passthrough{\lstinline!covered\_y * Δx!}。坐标压缩要保留相邻坐标差，节点代表小区间而非点。

\begin{lstlisting}[language={C++}]
struct Event { long long x, y1, y2; int delta; };
sort(events.begin(), events.end(), [](auto a, auto b) { return a.x < b.x; });
long long area = 0, last_x = events[0].x;
for (auto e : events) {
    area += covered_length() * (e.x - last_x);
    update(y1_id(e.y1), y2_id(e.y2) - 1, e.delta);
    last_x = e.x;
}
\end{lstlisting}

\section{最近点对}\label{ux6700ux8fd1ux70b9ux5bf9}

分治按 x 排序，递归求左右最小距离，只检查中线附近按 y 排序的常数个点，复杂度 \passthrough{\lstinline!O(n log n)!}。若只需竞赛常见规模，网格哈希也可用，但要证明桶大小与最小距离的关系。

\begin{lstlisting}[language={C++}]
long double closest(vector<Point>& p, int l, int r) {
    if (r - l <= 3) {
        long double ans = 1e100L;
        for (int i = l; i < r; ++i) for (int j = l; j < i; ++j) ans = min(ans, dist2(p[i], p[j]));
        sort(p.begin() + l, p.begin() + r, [](Point a, Point b) { return a.y < b.y; });
        return ans;
    }
    int m = (l + r) / 2; Real mid = p[m].x;
    Real ans = min(closest(p, l, m), closest(p, m, r));
    vector<Point> strip;
    for (int i = l; i < r; ++i) if ((p[i].x - mid) * (p[i].x - mid) < ans) strip.push_back(p[i]);
    for (int i = 0; i < (int)strip.size(); ++i)
        for (int j = i + 1; j < (int)strip.size() && strip[j].y - strip[i].y < sqrtl(ans); ++j)
            ans = min(ans, dist2(strip[i], strip[j]));
    return ans;
}
\end{lstlisting}

\section{半平面交（Half-Plane Intersection，HPI，了解）}\label{ux534aux5e73ux9762ux4ea4half-plane-intersectionhpiux4e86ux89e3}

把每个半平面写成有向直线左侧，按极角排序并用双端队列维护交集；处理平行情形和角度相等是难点。若题目只问凸多边形裁剪，优先用逐边 Sutherland--Hodgman。

\begin{lstlisting}[language={C++}]
vector<Point> clip(const vector<Point>& poly, Point a, Point b) {
    // 保留有向边 a->b 的左侧；cp/cq 是端点相对切线的叉积符号。
    vector<Point> out;
    for (int i = 0; i < (int)poly.size(); ++i) {
        Point p = poly[i], q = poly[(i + 1) % poly.size()];
        Real cp = cross(b - a, p - a), cq = cross(b - a, q - a);
        // p 在合法侧就保留 p。
        if (cp >= -EPS) out.push_back(p);
        // 一进一出时，边 pq 与切线有交点，按线性插值求交。
        if ((cp >= 0) != (cq >= 0)) {
            Real t = cp / (cp - cq);
            out.push_back(p + (q - p) * t);
        }
    }
    return out;
}
\end{lstlisting}

下面是``有向直线左侧为合法半平面''的真正半平面交模板；返回空集表示交集无界或为空时需结合题目另行处理，常规竞赛题通常会额外给一个大矩形作边界。

\begin{lstlisting}[language={C++}]
struct HalfPlane { Point p, v; Real ang; };
Point line_inter(HalfPlane a, HalfPlane b) {
    // a.p+t*a.v 与 b.p+s*b.v 的交点参数。
    Real t = cross(b.p - a.p, b.v) / cross(a.v, b.v);
    return a.p + a.v * t;
}
bool outside(HalfPlane h, Point p) {
    // 有向直线左侧合法；叉积小于 0 表示点在右侧，应该删掉。
    return sgn(cross(h.v, p - h.p)) < 0;
}

vector<Point> half_plane_intersection(vector<HalfPlane> h) {
    // 极角排序让平面边界按方向排列，双端队列只需维护首尾交点。
    for (auto &x : h) x.ang = atan2l(x.v.y, x.v.x);
    sort(h.begin(), h.end(), [](const HalfPlane& a, const HalfPlane& b) {
        return a.ang < b.ang;
    });
    deque<HalfPlane> q;
    for (auto cur : h) {
        // 新半平面会淘汰队尾/队首不在新半平面内的交点。
        while (q.size() >= 2 && outside(cur, line_inter(q[q.size() - 2], q.back()))) q.pop_back();
        while (q.size() >= 2 && outside(cur, line_inter(q[0], q[1]))) q.pop_front();
        if (!q.empty() && sgn(cross(q.back().v, cur.v)) == 0) {
            // 同向平行线只保留更严格的一条，反向平行通常意味着交集为空。
            if (outside(cur, q.back().p)) q.back() = cur;
        } else q.push_back(cur);
    }
    while (q.size() >= 3 && outside(q.front(), line_inter(q[q.size() - 2], q.back()))) q.pop_back();
    while (q.size() >= 3 && outside(q.back(), line_inter(q[0], q[1]))) q.pop_front();
    if (q.size() < 3) return {};
    vector<Point> poly;
    for (int i = 0; i < (int)q.size(); ++i)
        poly.push_back(line_inter(q[i], q[(i + 1) % q.size()]));
    return poly;
}
\end{lstlisting}

\section{几何中的随机化与哈希技巧}\label{large-geometry-random-hash}

\begin{itemize}
\tightlist
\item
  随机 \passthrough{\lstinline!shuffle!} + 快速选择求期望线性第 k 小。
\item
  64 位随机哈希表示集合，支持多重集合差异的高概率比较；严谨题目用排序/Trie/计数。
\item
  Treap 随机优先级避免退化；不要把 \passthrough{\lstinline!rand()!} 作为高质量随机源，使用 \passthrough{\lstinline!mt19937\_64!}。
\end{itemize}

\begin{lstlisting}[language={C++}]
mt19937_64 rng(chrono::steady_clock::now().time_since_epoch().count());
shuffle(a.begin(), a.end(), rng);
nth_element(a.begin(), a.begin() + k, a.end());
int answer = a[k];
uint64_t random_hash = rng();
\end{lstlisting}

\section{几何中的位集优化}\label{large-geometry-bitset}

\passthrough{\lstinline!std::bitset!} 将 64/机器字并行，适合传递闭包（布尔矩阵）、集合交并、子集计数；\passthrough{\lstinline!bitset!} 大小需编译期常量，动态场景使用 \passthrough{\lstinline!vector<unsigned long long>!}。

\begin{lstlisting}[language={C++}]
const int N = 2000;
bitset<N> reach[N];
for (int k = 0; k < n; ++k)
    for (int i = 0; i < n; ++i)
        if (reach[i][k]) reach[i] |= reach[k];
\end{lstlisting}

\section{几何快速小技巧}\label{large-geometry-quick-tips}

\begin{itemize}
\tightlist
\item
  \passthrough{\lstinline!std::gcd!}、\passthrough{\lstinline!std::lcm!}（C++17）。
\item
  \passthrough{\lstinline!\_\_builtin\_ctzll!} 只对非零数调用。
\item
  排序后用 \passthrough{\lstinline!unique!} 去重，差异数组用 \passthrough{\lstinline!adjacent\_difference!}。
\item
  大量模乘、叉积、斜率比较时先判断上界，必要时统一 \passthrough{\lstinline!\_\_int128!}。
\end{itemize}

\begin{lstlisting}[language={C++}]
long long g = std::gcd(a, b);
long long l = a / g * b;
sort(v.begin(), v.end());
v.erase(unique(v.begin(), v.end()), v.end());
long long lowbit = x & -x; // x != 0
long long safe_product = (long long)((__int128)x * y % MOD);
\end{lstlisting}

